\section{Liquidity Sources and the State Boundary}
\label{sec:sources}

We formalize the three sources by the only property that matters to the
router: what state their price is a function of, and therefore where
that price can be computed.

\subsection{Native AMM: price is a pure function of reserves}

A constant-function market maker prices a trade as a deterministic
function of its on-chain reserves $R$. For V2 the invariant is the
constant product $x \cdot y = k$; a swap of $\Delta_{\mathrm{in}}$ (net
of fee $\phi$) returns
\[
\Delta_{\mathrm{out}}
= \frac{y \,\cdot\, (1-\phi)\,\Delta_{\mathrm{in}}}
       {x + (1-\phi)\,\Delta_{\mathrm{in}}}.
\]
For V3 the same invariant holds piecewise over each initialized tick
range, with liquidity $L$ concentrated between price bounds; the
realized output is the integral of $L$ across the ticks the trade
crosses (\cite{uniswapv3}). The native constant-function AMM comprises
these V2 and V3 contracts only; the \texttt{0x9999} precompile is
\emph{not} a constant-function AMM but the DEX receipt-settlement conduit
for the D-Chain limit-order-book venue, whose price is set by the order
book rather than by on-chain reserves.

The decisive property is purity: given the block's reserves, the output
is exact and reproducible by any party. No off-chain information is
required to quote a native pool, so a native-pool quote is trustless
and needs no signature.

\subsection{Signed-quote market maker: price depends on off-chain state}

The PMM prices a pair $A \to B$ as
\[
y \;=\; \mathrm{ref}_{A \to B}\cdot(1-s)\cdot x,
\]
where $\mathrm{ref}_{A\to B}$ is an external reference price and $s$ is
the spread. Both inputs are off-chain: the reference comes from a
venue's order book, and the spread is set against inventory and
volatility. The quote therefore \emph{cannot} be a pure on-chain
function --- it is a signed assertion by the market maker, valid until a
stated expiry, and made enforceable on-chain only by verifying the
signature (\S\ref{sec:rfq}).

After an on-chain fill the PMM is left with an inventory imbalance ($-x$
of $A$, $+y$ of $B$ relative to target) and restocks by trading the
opposite direction on the external venue. The spread $s$ is the
compensation for the risk taken between quote and completed hedge
(\S\ref{sec:rfq-risk}).

\subsection{External DEX: trustless but last}

External on-chain pools are, like native AMMs, pure functions of their
reserves and require no signature. They differ only in policy: they are
consulted last, purely to convert a would-be revert into a fill, and at
whatever price they offer. They introduce no new trust assumption ---
the settlement transaction either obtains the required output from them
or reverts.

\subsection{Summary: the boundary as a table}

Table~\ref{tab:sources} states the property that drives every later
decision. The router (\S\ref{sec:router}) reads off this table; the
backend/no-backend split (Principle~\ref{prin:dividing}) is exactly the
``off-chain state?'' column.

\begin{table}[h]
\centering
\caption{The three sources by state dependence.}
\label{tab:sources}
\small
\begin{tabular}{@{}llll@{}}
\toprule
Source & Price is a function of & Off-chain state? & Needs signature? \\
\midrule
Native AMM V2/V3/V4 & on-chain reserves & no  & no \\
Signed-quote PMM    & reference + inventory + spread & yes & \textbf{yes} \\
External DEX        & on-chain reserves & no  & no \\
\bottomrule
\end{tabular}
\end{table}
